% Generated draft: requires mathematical and source review.
\documentclass[11pt]{article}
\usepackage[margin=25mm]{geometry}
\usepackage{fontspec}
\setmainfont{Latin Modern Roman}
\usepackage{amsmath,amssymb,mathtools,unicode-math}
\setmathfont{Latin Modern Math}
\usepackage{enumitem,needspace,xurl,hyperref}
\setlength{\parindent}{0pt}
\setlength{\parskip}{5pt}
\setlength{\emergencystretch}{3em}
\newcommand{\N}{\mathbb{N}}
\newcommand{\Z}{\mathbb{Z}}
\newcommand{\Q}{\mathbb{Q}}
\newcommand{\R}{\mathbb{R}}
\newcommand{\C}{\mathbb{C}}
\newcommand{\F}{\mathbb{F}}
\newcommand{\A}{\mathbb{A}}
\newcommand{\PP}{\mathbb{P}}
\newcommand{\E}{\mathbb{E}}
\newcommand{\eps}{\varepsilon}
\newcommand{\dd}{\,\mathrm{d}}
\begin{document}
\section*{UnsolvedMath ID: 30003892}
\textbf{Source problem code:} catalogue entry
\subsection{Definitions and mathematical statement}
Let
\[
S^3=\{(z_1,z_2)\in\C^2:|z_1|^2+|z_2|^2=1\}.
\]
A smooth link in \(S^3\) is a smooth embedded finite disjoint union of circles. Two such links have the same link type if they are ambient isotopic in \(S^3\). Whether orientations are part of the link type, and whether the empty link is allowed, are not specified in the supplied record.

Under the usual interpretation of an affine complex plane algebraic curve, write
\[
V_f=\{(z_1,z_2)\in\C^2:f(z_1,z_2)=0\},
\qquad f\in\C[z_1,z_2]\setminus\C.
\]
Here the underlying zero set is used, so polynomial multiplicities do not change the intersection. For a smooth link \(L\subset S^3\), define the realizability condition
\[
\mathcal C(L)\quad\Longleftrightarrow\quad
\begin{gathered}
\text{there exist }f\in\C[z_1,z_2]\setminus\C
\text{ and a smooth ambient isotopy }H_t:S^3\longrightarrow S^3,\\
0\leq t\leq1,\qquad H_0=\operatorname{id}_{S^3},\qquad
H_1(L)=V_f\cap S^3.
\end{gathered}
\]
The question is to characterize the link types satisfying \(\mathcal C(L)\), and to identify topological obstructions to this realizability condition. A complete characterization would give a condition on the link type that is necessary and sufficient for \(\mathcal C(L)\).

This formulation makes explicit one interpretation of the extracted question; its intended hypotheses require review. The source does not require \(V_f\) to be irreducible or nonsingular, nor does it state that the intersection with \(S^3\) is transverse. Equality with \(H_1(L)\) above requires the intersection itself to be a smooth link, but does not impose transversality. If the intended definition of a \(\C\)-link requires \(V_f\) to be nonsingular near the sphere and transverse to it, those are additional hypotheses defining a potentially different target class. In that setting the complex orientation of the curve and the boundary orientation of its portion inside the unit ball supply an orientation on the link; the source does not say whether this orientation must be respected. The sphere is the stated unit sphere, with no supplied requirement that it be a sufficiently small sphere around an isolated curve singularity.
\subsection{Short English statement}
Determine exactly which link types in the unit \(3\)-sphere occur as intersections with complex algebraic curves in \(\C^2\), and find topological criteria or obstructions for such realizations. The intended smoothness, transversality, and orientation conventions remain unresolved.
\subsection{Sources}
\begin{itemize}
\item Supplied local UnsolvedMath registry record: permanent ID \(30003892\), title \emph{Topological Characterization of Complex Algebraic Links}, exact source problem code \texttt{catalogue entry}, in the Oberwolfach Reports Open Problems collection. No individual proposer is supplied; the proposed year is recorded as 2018.
\item Supplied report citation: \emph{Topologie (2019)}, Oberwolfach Reports, pp.\ 1857--1910, DOI \texttt{10.4171/owr/2018/31}. \url{https://doi.org/10.4171/owr/2018/31}. The record attributes the question to workshop open-problem discussions but supplies no problem-specific page or speaker attribution.
\item Supplied literature source: Borodzik--Hedden--Livingston, plane algebraic curves. \url{https://arxiv.org/abs/1409.2111}.
\item Supplied literature source: Oberwolfach report. \url{https://doi.org/10.4171/owr/2018/31}.
\item Literature assessment checked at \texttt{2026-08-21}: the supplied curation labels the problem open and states that known Heegaard--Floer and algebraic-geometric obstructions address cases without settling the general question. This is supplied dated curation, not an independently verified assessment. The background expressly requests independent verification of time-sensitive status and cited claims.
\item The registry labels the statement \texttt{corrected\_verified} and says that notation, typography, and omitted definitions were repaired while preserving mathematical content. This is a supplied editorial qualification, not independent mathematical verification of this draft.
\end{itemize}
\end{document}
